\documentclass[12pt]{article}

% Layout.
\usepackage[top=1in, bottom=0.75in, left=0.75in, right=0.75in, headheight=1in, headsep=6pt]{geometry}

% Fonts.
\usepackage{mathptmx}

\usepackage[scaled=0.86]{helvet}
\renewcommand{\emph}[1]{\textsf{\textbf{#1}}}

% Misc packages.
\usepackage{amsmath,amssymb,latexsym,amsthm}
\usepackage{graphicx,hyperref}
\usepackage{array}
\usepackage{xcolor}
\usepackage{multicol}
\usepackage{tabularx,colortbl}
\usepackage{enumitem}
\usepackage{soul,multicol}
\usepackage{setspace}


\hypersetup{
    colorlinks=true,
    linkcolor=blue,
    filecolor=magenta,      
    urlcolor=blue,
    pdftitle={Proofs Worksheet},
    pdfpagemode=FullScreen,
    }

\def\mailto#1{\href{mailto:#1}{#1}}

%% special math symbols
\newcommand{\N}{\mathbb{N}}
\newcommand{\Z}{\mathbb{Z}}
\newcommand{\R}{\mathbb{R}}
\newcommand{\Q}{\mathbb{Q}}
\newcommand{\sse}{\subseteq}
\newcommand{\mcp}{\mathcal{P}}
\newcommand{\ol}[1]{\overline{#1}}

%other specials
\newcommand{\be}{\begin{enumerate}}
\newcommand{\ee}{\end{enumerate}}
\newcommand{\se}{\subseteq}
\newcommand{\spe}{\supseteq}
\newcommand{\ds}{\displaystyle}
\newcommand{\lcm}{\text{lcm}}
%\newcommand{\gcd}{\text{gcd}}

% Paragraph spacing
\parindent 0pt
\parskip 6pt plus 1pt
\def\tableindent{\hskip 0.5 in}
\def\ts{\hskip 1.5 em}

%header
\usepackage{fancyhdr}
\pagestyle{fancy} 
\lhead{\large\sf\textbf{MATH 265: Introduction to Mathematical Proofs}}
\rhead{\large\sf\textbf{Worksheet 21: Ch 11/12}}

\newcommand{\localhead}[1]{\par\smallskip\textbf{#1}\nobreak\\}%
\def\heading#1{\localhead{\large\emph{#1}}}
\def\subheading#1{\localhead{\emph{#1}}}

\newenvironment{clist}%
{\bgroup\parskip 0pt\begin{list}{$\bullet$}{\partopsep 4pt\topsep 0pt\itemsep -2pt}}%
{\end{list}\egroup}%

\begin{document}
\begin{enumerate}
\item Quick Review
\vspace*{-0.5cm}
	\begin{enumerate}
	\item $R$ is an equivalence relation on $\R$ defined as $xRy$ if $\sin(x)=\sin(y).$ Describe and/or draw its equivalence classes.
	\vfill 
	
	\item We say $S$ is a partition of the set $A$ if 	
	\vspace{.75in}
	\end{enumerate}
\item Theorem 11.1: Suppose $R$ is an equivalence relation on the set $A.$ For every $a, b \in A,$ $[a]=[b]$ if and only if $aRb.$
\vfill
\item Theorem 11.2: Suppose $R$ is an equivalence relation on the set $A$. The set $S=\{[a] : a \in A\}$ is a partition of $A.$
\vfill
\newpage
\item Definition: The \textbf{integers modulo $n$}, denoted $\Z_n$, is the set $\{[0],[1], \cdots, [n-1]\}$ of equivalences classes of $\Z$ modulo $n.$\\
Try performing arithmetic in $\Z_{10}$. Does it depend on the representative of the equivalence class?
	\begin{multicols}{3}
		\begin{enumerate}
		\item $[0]+[6]=$\\
		\item $[2]+[3]= $\\
		\item $[9]+[5]=$\\
		\item $[3] - [7]=$\\
		\item $[6] \cdot [1]=$\\
		\item $[3] \cdot [7]=$	\\
		\item $[4] \cdot [5]=$	\\
		\item $[3] \div [7]=$\\	
		\item $[3] \div [4]=$\\
		\end{enumerate}
	\end{multicols}
	\vfill
\item What do you \emph{think} should be the definition of a relation $R$ from the set $A$ to the set $B$?
\vfill	
\item Think up a relation $R$ from $A=\{a,b\}$ to $B=\{1,2,3\}$. How many distinct relations from $A$ to $B$ are possible?
\vfill
\item How many relations $R$ from $A$ to $B$ are functions? (This is require you to decide what you think should be the definition of a \emph{function} from $A$ to $B$...)
\vfill
\item Think up a relation $S$ from $C=\{a,b,c\}$ to $D=\{1,2,3,4,5\}$ that is a function. Determine the domain and range of $S$.
\vfill
\item In the context of relations that are functions where we already have ideas about the words \emph{domain} and \emph{range}, what might the term \emph{codomain} mean?
\vfill
\end{enumerate}
\end{document}










